\section{Decision-Support Use in Scientific Review}
\label{sec:decision-support-use}

ClaimBench is intended to be used before and during manuscript review. The starting point is a scientific AI study that already contains computational artifacts and a draft interpretation. The decision-support task is then to determine which statements can remain in the paper, which ones must be weakened, and which ones require additional experiments. This is a practical problem in research management because reviewers rarely object only to a number. They object to the interpretation attached to that number.

The workflow begins with claim discovery. Candidate statements are extracted from the manuscript, README files, and other public-facing project documents. In the reproducible configuration, this extraction is deterministic. If an external LLM is used, its output is loaded as a local candidate list and is not treated as evidence. This distinction is important because it keeps the human writing workflow separate from the scientific authorization workflow. A generated sentence can be useful for inspection and still be unsupported.

The second step is claim classification. Each candidate statement is mapped to a claim family. A statement such as ``the model predicts held-out responses'' belongs to a predictive family. A statement such as ``the model identifies a mechanism'' belongs to a stronger mechanistic family. A statement such as ``the system is a digital twin'' belongs to an even more demanding scope family. The classification determines which evidence dimensions are required. It also determines which weaker wording remains valid when stronger evidence is absent.

The third step is artifact binding. Each claim must be linked to an executable or stored artifact. For synthetic cases, the artifact contains the known-truth configuration and the gate decision. For SciFact, it contains retrieval, rationale, abstention, and overclaiming-risk information. For Tuebingen, it contains direction attempts and causal-overclaiming controls. For release reproducibility, it contains the required artifact list, the expected decisions, and the dirty-state check. This binding prevents a manuscript from making a claim that is no longer connected to the code and data used to justify it.

The fourth step is decision generation. ClaimBench does not return a single global acceptance score. It returns a claim-level status. A supported claim can remain if its wording matches the evidence boundary. A blocked claim must be removed or downgraded. An uncertain claim can be reported only as unstable, exploratory, or requiring confirmation. This makes the output suitable for author revision because it tells the author what action is needed, not only that a problem exists.

\begin{table}[t]
\centering
\caption{Decision states returned by ClaimBench and their intended editorial consequence. The output is not a ranking of models, but a concrete action on manuscript wording.}
\label{tab:decision-actions}
\small
\begin{tabularx}{\textwidth}{p{0.18\textwidth}p{0.28\textwidth}X}
\toprule
Decision state & Evidence status & Recommended manuscript action \\
\midrule
Supported & All required evidence blocks are present and stable under the declared controls & Keep the claim, but preserve the bounded wording and cite the decision artifact \\
Blocked & At least one non-interchangeable evidence block is missing & Remove the claim or replace it with a weaker statement supported by the available evidence \\
Uncertain & Evidence is present but unstable, threshold-sensitive, or insufficiently confirmed & Report the result as exploratory or requiring confirmation, not as established support \\
Out of scope & The claim family requires evidence not collected by the study & State the limitation explicitly and avoid implying that the framework has evaluated it \\
Needs external review & The decision depends on domain judgement outside the executable artifact & Use the artifact as review support, but keep the final judgement with human experts \\
\bottomrule
\end{tabularx}
\end{table}


This process also supports reviewer communication. If a reviewer argues that a mechanistic claim is unsupported, the authors can point to the corresponding claim contract and show whether topology, direction, and alternative-control evidence were present. If a reviewer argues that the framework is only a collection of rules, the authors can point to the ablation artifact showing which risks reappear when each rule is removed. If a reviewer argues that LLMs have introduced an opaque decision layer, the authors can point to the LLM audit showing that no API was called and that no LLM candidate was authoritative.

The system is therefore not a substitute for reviewer judgement. It is a structured way of making reviewer judgement easier to apply. It reduces ambiguity by forcing each high-level statement to carry its evidential requirements. It also reduces post-hoc expansion because the release package records which claims were allowed at the time of the analysis. This is particularly useful for applied AI papers, where it is common to present a technically correct result and then move too quickly toward stronger scientific or operational interpretation.

In practice, ClaimBench can be used in three modes. In author-preparation mode, it acts as a pre-submission audit that identifies overclaiming before the paper is sent to a journal. In internal-review mode, it allows a research group to compare alternative manuscript framings against the same evidence package. In reviewer-response mode, it provides a traceable structure for responding to criticisms about evidence, causality, reproducibility, or scope. These modes are aligned with the DSS framing because the system supports a sequence of decisions rather than a single automated judgement.

The current manuscript uses ClaimBench in all three senses. The package audits its own claim boundaries, reports external controls that stress the gate, and includes ablations that show why the components are needed. This self-application does not prove that the framework will generalize to every scientific AI domain. It does show that the proposed decision-support structure is executable and that it can be inspected by a reader without relying on informal assurances.

